\documentclass{article}
\usepackage[T1]{fontenc}
\usepackage{lmodern}
\usepackage{graphicx}
\usepackage{amsmath,amssymb}
\usepackage[a4paper,margin=2cm]{geometry}
\usepackage[absolute,overlay]{textpos}
\setlength{\TPHorizModule}{1mm}
\setlength{\TPVertModule}{1mm}
\usepackage[indonesian]{babel}
\usepackage{hyperref}
\setlength{\emergencystretch}{2em}
% unit: Halaman 14

\begin{document}

% === Halaman 14 (python/native) ===

14

Enkripsi: 

1010  0010  0011  1010  1001 

1011  1011  1011  1011  1011  


Hasil XOR:  


0001  1001  1000  0001  0010     

Geser 1 bit ke kiri: 

0010  0011  0001  0010  0100  

Dalam notasi HEX: 

   2    3 

  1     2     4 

Jadi, hasil enkripsi plainteks  

10100010001110101001 

(A23A9 dalam notasi HEX) 


adalah 

00100011000100100100 

(23124 dalam notasi HEX) 

EK(P) = (P  K) << 1

\end{document}
